\documentclass[11pt,a4paper]{article}
\usepackage[margin=2.5cm]{geometry}
\usepackage{amsmath}
\usepackage[T1]{fontenc}
\usepackage[utf8]{inputenc}

\title{Method Explanation -- Group 20\\
\large Sustainable MRP 90-Minute Challenge}
\author{Group 20}
\date{}

\begin{document}
\maketitle

\section*{Forecast}
We used multiplicative seasonal decomposition ($Y = T \cdot S \cdot R$).
We computed monthly seasonal indices from 104 weeks of history, then removed seasonality
from demand and fitted a linear model (\texttt{fitlm}) using \textit{WorkingDays},
\textit{ConfirmedOrders} and \textit{EcoTenderWeek} as predictors (other variables dropped
by backward $p$-value selection to avoid overfitting). Validation with Leave-One-Out
jackknife gave MAPD $= 0.10\%$ (very good). Safety stock $= 1.65 \times \sigma_{\text{LOO}} \approx 2$ units.

\section*{Production}
We produce at maximum weekly capacity every week. Stockout and service-level penalties
together amount to 430~EUR/unit, far above production cost (42~EUR/unit). Since total
capacity is below total forecast, producing less would only increase losses.

\section*{Procurement}
We solved a linear program (\texttt{linprog}) to minimise purchase cost plus material
holding cost, subject to: material balance with deterministic lead times, supplier weekly
capacity caps, carbon cap $\leq 155{,}000$~kgCO$_2$ (hard constraint), and low-carbon
mass share $\geq 35\%$ (hard constraint). Hard constraints were chosen because penalties
are much higher than the marginal cost of complying. Result: zero sustainability penalties.

\end{document}
